\documentclass[11pt]{article}
\usepackage[margin=1in]{geometry}
\usepackage{amsmath,amssymb}
\usepackage[T1]{fontenc}
\usepackage{lmodern}
\usepackage{microtype}
\usepackage{xurl}
\usepackage{booktabs}
\usepackage{graphicx}
\usepackage{hyperref}
\setlength{\emergencystretch}{3em}

\title{Fertility Project Experiment Note: Model Setup Trial + Empirical Idea Ranking}
\author{Project 03: Fertility and Housing Supply}
\date{February 20, 2026}

\begin{document}
\maketitle

\section*{Purpose}
This note does two things:
\begin{enumerate}
\item Runs a small computational experiment based on the model setup in
\path{projects/03_fertility_and_housing_supply/notes/model/02_fertility_extension_model_key_equations.tex}.
\item Ranks empirical causal ideas and expands the top 2 in detail (including a NIMBY-focused design).
\end{enumerate}

\section{Part I: Model setup trial (code + run)}

\subsection*{What was implemented}
A version-0 transition simulation was coded in:
\begin{itemize}
\item \path{projects/03_fertility_and_housing_supply/code/fertility_extension_experiment.py}
\end{itemize}
The prototype keeps the model intentionally stylized but consistent with the note's key feedback loop:
\[
\text{prices} \rightarrow \text{fertility} \rightarrow \text{cohort shares} \rightarrow \text{political tightness} \rightarrow \text{prices}.
\]

Scenarios simulated for 50 periods:
\begin{itemize}
\item No damping (aggressive update).
\item Damping 0.25 (stabilized fixed-point style update).
\item Damping 0.25 + family-sized supply policy shift from period 15 onward.
\end{itemize}

Outputs were written to:
\begin{itemize}
\item \path{projects/03_fertility_and_housing_supply/notes/build/model_experiment_summary.csv}
\item \path{projects/03_fertility_and_housing_supply/notes/build/model_experiment_paths.png}
\end{itemize}

\subsection*{Run command}
\begin{verbatim}
python3 projects/03_fertility_and_housing_supply/code/fertility_extension_experiment.py
\end{verbatim}

\subsection*{Results (from summary csv)}
\begin{center}
\begin{tabular}{lccc}
\toprule
Scenario & Avg price (t=10..49) & Avg fertility (t=10..49) & Max abs price change \\
\midrule
No damping & -5.939 & 0.723 & 2.426 \\
Damping 0.25 & -2.597 & 0.416 & 1.948 \\
Damping 0.25 + policy & -2.638 & 0.443 & 1.349 \\
\bottomrule
\end{tabular}
\end{center}

\textbf{Interpretation:}
\begin{itemize}
\item The undamped update is much less stable (larger jumps and corner-like dynamics).
\item Damping materially improves numerical behavior.
\item The family-sized supply shift raises fertility in the damped run (0.416 to 0.443 average fertility rate in this toy environment).
\end{itemize}

\textbf{What failed and what was changed:}
Initial runs generated overflow in exponential terms (numerical blow-up), so the code was patched with bounded exponent arguments and price clamps. This is direct evidence that the full model will require strong numerical safeguards (damping, normalization, bounded controls).

\begin{figure}[h!]
\centering
\includegraphics[width=0.9\textwidth]{build/model_experiment_paths.png}
\caption{Version-0 paths: damping sensitivity and policy experiment}
\end{figure}

\subsection*{Value of this experiment}
This is not structural estimation. It is a feasibility stress test that confirms the convergence warning in the model note is operationally relevant, not just theoretical.

\section{Part II: Empirical causal ideas ranked}

\subsection*{Ranking criteria}
Scores reflect four dimensions: identification credibility, data feasibility, novelty potential, and alignment with your housing-fertility channel.

\begin{center}
\begin{tabular}{p{0.44\textwidth}c}
\toprule
Idea & Rank \\
\midrule
State middle-housing legalization (staggered reform) and fertility & 1 \\
NIMBY political-economy design: close local elections and pro-housing wins & 2 \\
Family-sized unit-composition shock and fertility (shift-share IV) & 3 \\
Short-term rental crackdown and fertility & 4 \\
Single-city zoning mega-reform with legal interruptions & 5 \\
\bottomrule
\end{tabular}
\end{center}

\section{Top idea 1 (expanded): State middle-housing legalization and fertility}
\textbf{Question:} Do legalizations of duplex/triplex/fourplex construction increase births among prime childbearing cohorts?

\textbf{Estimand:} Dynamic ATT of reform exposure on age-specific birth rates.

\textbf{Design:} Staggered DiD/event study at county or PUMA level with intensity-by-exposure.

\textbf{Core specification (event study):}
\[
BirthRate_{it} = \sum_{k \neq -1}\beta_k \cdot 1\{EventTime_{it}=k\} + \alpha_i + \delta_t + X_{it}'\Gamma + \varepsilon_{it}.
\]
Intensity extension:
\[
BirthRate_{it} = \beta\,(Post_{it}\times Exposure_i) + \alpha_i + \delta_t + X_{it}'\Gamma + \varepsilon_{it}.
\]

\textbf{Data (feasible now):}
\begin{itemize}
\item CDC WONDER natality (county-level births by mother's age).
\item Census Building Permits Survey (new units; single vs multifamily).
\item FHFA/Zillow house-price and rent controls.
\item Reform timing panel (state/city policy adoption and implementation).
\end{itemize}

\textbf{Why this is strong:}
\begin{itemize}
\item Direct policy shock on the supply margin.
\item Directly matches your mechanism (housing constraints $\\rightarrow$ family formation/fertility).
\item Good external validity due to multi-jurisdiction rollout.
\end{itemize}

\textbf{Main threats and tests:}
\begin{itemize}
\item Threat: pre-trends in treated areas.
\item Test: pre-period event coefficients and placebo reform dates.
\item Threat: migration composition shifts.
\item Test: control for migration proxies; run residence-duration restricted samples if possible.
\item Threat: weak first stage on effective supply.
\item Test: document permit/unit-mix response as first-stage evidence.
\end{itemize}

\textbf{Novelty claim (current confidence): Medium-High.}
The novelty is to estimate fertility effects of explicit middle-housing legalizations, not only generic price shocks.

\section{Top idea 2 (expanded): NIMBY political-economy design via close elections}
\textbf{Question:} Does narrowly electing pro-housing local governments causally increase supply and, through this channel, affect fertility?

\textbf{Estimand:}
\begin{itemize}
\item First stage: effect of pro-housing win on permits/unit composition and prices.
\item Second stage: effect of induced housing-supply changes on fertility.
\end{itemize}

\textbf{Design:} Close-election RD at municipal level + 2SLS mediation.

First stage:
\[
Supply_{it} = \pi\,ProHousingWin_{it} + f(Margin_{it}) + \eta_i + \tau_t + u_{it}.
\]
Second stage:
\[
BirthRate_{it} = \beta\,\widehat{Supply}_{it} + f(Margin_{it}) + \eta_i + \tau_t + e_{it}.
\]

\textbf{Why this is NIMBY-relevant and novel:}
\begin{itemize}
\item It directly uses political turnover in the local anti-/pro-development margin.
\item It maps your theoretical political channel into a causal empirical design.
\item Compared with standard housing papers, the fertility endpoint is a newer margin.
\end{itemize}

\textbf{Data requirements (hard but feasible):}
\begin{itemize}
\item Municipal election returns with candidate coding (pro-housing stance).
\item Municipal permits and unit mix.
\item Local prices/rents.
\item Birth outcomes at the finest feasible geography (county minimum; subcounty where available).
\end{itemize}

\textbf{Main risks:}
\begin{itemize}
\item Candidate stance coding may be noisy.
\item Fertility outcomes may be coarser than election geography.
\item Sample size can be thin once restricting to close races.
\end{itemize}

\textbf{Risk mitigation:}
\begin{itemize}
\item Use transparent text-based coding protocol with inter-rater checks.
\item Aggregate to county-year with municipality exposure weights.
\item Pool multiple states and cycles.
\end{itemize}

\textbf{Novelty claim (current confidence): Medium.}
Potentially high upside, but novelty depends on whether recent housing-politics papers already run near-identical RD-to-fertility pipelines.

\section{Pragmatic recommendation}
If the objective is ``credible result in finite time,'' start with Top 1.
If the objective is ``high-upside NIMBY contribution,'' develop Top 2 in parallel as a second track.

\section*{Deliverables generated in this session}
\begin{itemize}
\item Prototype code: \path{projects/03_fertility_and_housing_supply/code/fertility_extension_experiment.py}
\item Experiment outputs: \path{projects/03_fertility_and_housing_supply/notes/build/model_experiment_*}
\item This report PDF/TeX.
\end{itemize}

\end{document}
